\section{Dekooderipohjainen Transformer-arkkitehtuuri suurille
kielimalleille}

Luvuissa \ref{sec:discrete} ja \ref{sec:mlp} kielimallinnusta
tarkasteltiin seuraavan tokenin ehdollisen todennäköisyyden estimointina.
Olkoon \(V\) tokenisanasto ja \(\Delta(V)\) sanaston yli määriteltyjen
todennäköisyysjakaumien joukko. Kielimalli voidaan tällöin kuvata
funktiona, joka muuntaa aiemmin havaitun tokenikontekstin seuraavan tokenin
todennäköisyysjakaumaksi.

\(n\)-gram-mallissa kontekstin pituus on ennalta määrätty. Jos mallin
käyttämän kontekstin pituus on \(n\), niin se voidaan kuvata funktiona
\[
  f_{\mathrm{ngram}} \colon V^n \to \Delta(V).
\]

Myös luvussa \ref{sec:mlp} tarkasteltu MLP-pohjainen kielimalli käyttää
kiinteän pituista kontekstia, vaikka sen parametrisoitu esitystapa
mahdollistaa frekvenssipohjaista \(n\)-gram-mallia paremman
yleistämisen.

Nykyaikaiset suuret kielimallit
(engl.\ \emph{large language models}, LLM) perustuvat samaan
autoregressiiviseen ennustetehtävään, mutta niiden käsittelemän kontekstin
pituus ei ole kiinteä. Olkoon \(L\)
mallin suurin sallittu kontekstinpituus. Dekooderipohjainen
Transformer-kielimalli voidaan kuvata funktiona
\[
  f_{\mathrm{Tr}} \colon \bigcup_{m=0}^{L} V^m \to \Delta(V).
\]
Malli voi siten muodostaa seuraavan tokenin jakauman millä tahansa
enintään \(L\) tokenin pituisella kontekstilla. Käytännössä suurin
kontekstinpituus on edelleen arkkitehtuurin ja laskentaresurssien
rajoittama, mutta mallin parametrien ei tarvitse olla sidottuja yhteen
tiettyyn kontekstin pituuteen.

Tässä luvussa tarkastellaan erityisesti dekooderipohjaista
Transformer-arkkitehtuuria, joka muodostaa useiden generatiivisten
kielimallien, kuten GPT-malliperheen, matemaattisen perustan.
Transformer esiteltiin alun perin Vaswanin ym.\ työssä
\emph{Attention Is All You Need} vuonna 2017
\cite{vaswani2017attention}. Alkuperäinen arkkitehtuuri oli
konekäännökseen suunniteltu enkooderi--dekooderi-rakenne. Pian tämän jälkeen esitettiin, että pelkkään Transformer-dekooderiin perustuva autoregressiivinen kielimalli voidaan esikouluttaa suurilla tekstiaineistoilla ja hienosäätää useisiin kielenkäsittelytehtäviin \cite{radford2018improving}.

\subsection{Itsehuomio}

Transformer-mallin keskeinen mekanismi on itsehuomio (engl.\ \emph{self-attention}). Sen avulla tokenin esitys voidaan muodostaa koko syötekontekstin perusteella sen sijaan, että huomioitaisiin vain kiinteä määrä viimeisimpiä tokeneita. Kontekstin tokenien merkitys ei määräydy pelkän etäisyyden mukaan, vaan itsehuomio mahdollistaa kunkin tokenin kannalta olennaisen tiedon dynaamisen painottamisen yli koko sekvenssin.

Olkoon tarkasteltavan tokenijonon pituus \(T\) ja olkoon \(h_i \in \mathbb{R}^{d}\) position \(i\) tokenia vastaava esitysvektori. Näistä muodostetaan syötematriisi
\begin{align*}
  H
  =
  \begin{bmatrix}
    h_1^{\top} \\
    h_2^{\top} \\
    \vdots \\
    h_T^{\top}
  \end{bmatrix}
  \in \mathbb{R}^{T \times d}.
\end{align*}

Itsehuomiomekanismi muodostaa syöte-esityksistä kolme erillistä lineaarista projektiota: kyselyn (engl.\ \emph{query}), avaimen (engl.\ \emph{key}) ja arvon (engl.\ \emph{value}). Olkoot \(W_Q, W_K \in \mathbb{R}^{d \times d_k}\) ja \(W_V \in \mathbb{R}^{d \times d_v}\) opittavia painomatriiseja. Kysely-, avain- ja arvomatriisit lasketaan kertomalla syötematriisi näillä painoilla:
\begin{align*}
  Q &= H W_Q \in \mathbb{R}^{T \times d_k}, \\
  K &= H W_K \in \mathbb{R}^{T \times d_k}, \\
  V &= H W_V \in \mathbb{R}^{T \times d_v}.
\end{align*}
Intuitiivisesti kysely kuvaa sitä, millaista tietoa kukin tokeni hakee, avain kuvaa tokenin tarjoamaa informaatiota muille, ja arvo sisältää varsinaisen sisällön, jota yhdistetään uusiin esityksiin.

Kyselyiden ja avainten välinen vastaavuus lasketaan pistetuloina, joista muodostuu skaalattu huomiopistematriisi
\begin{align*}
  S = \frac{QK^{\top}}{\sqrt{d_k}} \in \mathbb{R}^{T \times T}.
\end{align*}
Matriisin alkio \(S_{ij}\) mittaa position \(i\) kyselyn ja position \(j\) avaimen yhteensopivuutta. Nimittäjä \(\sqrt{d_k}\) skaalaa pistetuloja, jotta niiden varianssi ei kasvaisi liian suureksi vektorien dimension mukana. Skaalaus ehkäisee softmax-funktion saturoitumista ja takaa stabiilimmat gradientit koulutuksen aikana \cite{vaswani2017attention}.

Huomiopisteet normalisoidaan todennäköisyysjakaumiksi soveltamalla softmax-funktiota riveittäin, jolloin saadaan huomiopainomatriisi
\begin{align*}
  A = \operatorname{softmax}(S) = \operatorname{softmax}\left(\frac{QK^{\top}}{\sqrt{d_k}}\right) \in \mathbb{R}^{T \times T},
\end{align*}
missä kukin alkio \(A_{ij} \in [0, 1]\) ilmaisee, kuinka suuri painoarvo position \(j\) arvolle annetaan position \(i\) uutta esitystä laskettaessa, ja riveille pätee \(\sum_{j=1}^{T} A_{ij} = 1\).

Lopullinen itsehuomion tulos lasketaan arvomatriisin rivien painotettuna summana matriisitulon avulla:
\begin{align*}
  \operatorname{Attention}(Q,K,V) = AV = \operatorname{softmax}\left(\frac{QK^{\top}}{\sqrt{d_k}}\right)V \in \mathbb{R}^{T \times d_v}.
\end{align*}
Tuloksena saadun matriisin kukin rivi \(i\) edustaa position \(i\) uutta kontekstualisoitua esitystä, johon on koottu tietoa kaikista sekvenssin tokeneista huomiopainojen \(A_{ij}\) mukaisesti.

Autoregressiivisissä kielimalleissa tavoitteena on mallintaa seuraavan tokenin todennäköisyyttä ja generoida tekstiä sekventiaalisesti, minkä vuoksi tokenin esitys ei saa hyödyntää tietoa tulevista tokeneista. Tätä varten itsehuomiota rajoitetaan \emph{kausaalisella maskauksella}, joka estää informaation vuotamisen tulevaisuudesta asettamalla position \(i\) huomiopainot nollaksi kaikille positioille \(j > i\). Tämä toteutetaan lisäämällä huomiopisteisiin ennen softmax-funktiota maskimatriisi \(M \in \mathbb{R}^{T \times T}\):
\begin{align*}
  M_{ij} =
  \begin{cases}
    0, & \text{kun } j \le i, \\
    -\infty, & \text{kun } j > i.
  \end{cases}
\end{align*}
Kausaalinen itsehuomio lasketaan tällöin muodossa
\begin{align*}
  \operatorname{CausalAttention}(Q,K,V) = \operatorname{softmax}\left(\frac{QK^{\top}}{\sqrt{d_k}} + M\right)V.
\end{align*}
Koska \(\exp(-\infty) = 0\), huomiopainomatriisista muodostuu alakolmiomatriisi. Tällöin kunkin tokenin esitys riippuu vain siitä itsestään ja sitä edeltävistä tokeneista, mikä säilyttää generointiprosessin kausaalisen rakenteen.

Yksittäisen huomiomekanismin ilmaisukykyä rajoittaa kuitenkin softmax-funktion luonne. Koska huomiopainomatriisille pätee \(A_{ij} \ge 0\) ja \(\sum_{j=1}^{T} A_{ij} = 1\), position \(i\) päivitetty esitysvektori \(o_i = \sum_{j=1}^{T} A_{ij} v_j\) on arvovektorien \(\{v_1, \dots, v_T\}\) konveksikombinaatio. Geometrisesti tämä tarkoittaa, että tulosvektori rajautuu arvovektorien virittämän konveksin verhon sisälle
\begin{align*}
  o_i \in \operatorname{conv}(\{v_1, \dots, v_T\}) = \left\{ \sum_{j=1}^{T} \alpha_j v_j \;\middle|\; \alpha_j \ge 0, \, \sum_{j=1}^{T} \alpha_j = 1 \right\}.
\end{align*}
Tämä asettaa huomiomekanismille geometrisen rajoitteen. Mikäli tokenin tulisi huomioida samanaikaisesti useita erilaisia kontekstuaalisia suhteita (kuten syntaktista rakennetta ja kaukaista semanttista viittausta), yksittäinen huomiopää joko keskittyy vain yhteen piirteeseen tai keskiarvoistaa useita piirteitä, jolloin erottelukyky heikkenee.

Tämän rajoitteen kiertämiseksi käytetään monipäistä itsehuomiota (engl.\ \emph{multi-head attention}), jossa huomiomekanismi jaetaan rinnakkaisiin esitysaliavaruuksiin \cite{vaswani2017attention}. Olkoon päiden lukumäärä \(h\). Kukin pää \(k \in \{1, \dots, h\}\) suorittaa itsenäisen itsehuomion omilla opittavilla projektioillaan:
\begin{align*}
  \operatorname{head}_k = \operatorname{Attention}(H W_k^Q, H W_k^K, H W_k^V),
\end{align*}
missä projektiomatriisit ovat \(W_k^Q, W_k^K \in \mathbb{R}^{d \times d_k}\) ja \(W_k^V \in \mathbb{R}^{d \times d_v}\). Tällöin kukin pää voi muodostaa oman konveksikombinaationsa toisistaan riippumatta omassa aliavaruudessaan.

Lopuksi eri päiden tuottamat esitykset ketjutetaan ja yhdistetään lineaarisella projektiolla takaisin mallin ulottuvuuteen \(d\):
\begin{align*}
  \operatorname{MultiHead}(H) = \left[ \operatorname{head}_1 \,\|\, \operatorname{head}_2 \,\|\, \dots \,\|\, \operatorname{head}_h \right] W^O,
\end{align*}
missä \(\|\) merkitsee sarakkeittaista ketjuttamista ja \(W^O \in \mathbb{R}^{h d_v \times d}\) on opittava parametrimatriisi. Lineaarinen yhdistäminen purkaa yksittäisten konveksien verhojen asettaman geometrisen rajoitteen ja mahdollistaa useiden erilaisten suhteiden samanaikaisen hyödyntämisen. Käytännössä pään dimensioiksi valitaan usein \(d_k = d_v = d / h\), jolloin monipäisen itsehuomion laskennallinen vaativuus vastaa karkeasti täyden dimension yksipäistä huomiota \cite{vaswani2017attention}.

\subsection{Positiokoodaus}
\label{sec:positional-encoding}

Itsehuomiomekanismi muodostaa tokenien väliset
huomiopisteet niiden esitysvektorien perusteella, mutta ei sellaisenaan
sisällä tietoa tokenien positioista. Ilman kausaalista maskia itsehuomio
on permutaatioekvivariantti eli syötteen tokenien permutoiminen
permutoi tulosteet samalla tavalla. Olkoon
\(P \in \{0,1\}^{T \times T}\) permutaatiomatriisi ja
\(H' = PH\) permutoitu syötematriisi. Tällöin kysely-, avain- ja
arvomatriiseille pätee

$$
\begin{aligned}
  Q' &= PHW_Q = PQ, \\
  K' &= PHW_K = PK, \\
  V' &= PHW_V = PV.
\end{aligned}
$$

Näin ollen permutoidut huomiopisteet ovat

$$
Q'K'^{\top}
=
(PQ)(PK)^{\top}
=
P QK^{\top}P^{\top}.
$$

Riveittäin laskettava softmax-funktio säilyttää tämän
permutaatiorakenteen, joten

$$
\operatorname{Attention}(PH)
=
P\operatorname{Attention}(H).
$$

Itsehuomio ei siis ilman erillistä positiotietoa kykene esittämään
tokenien järjestystä sisältöesityksistä riippumattomalla tavalla.
Luonnollisessa kielessä järjestys on kuitenkin merkityksen kannalta
olennainen. Samat tokenit eri järjestyksessä voivat muodostaa täysin
erilaisen ilmauksen.

Kausaalinen maski rikkoo yleisen permutaatioekvivarianssin, sillä se
määrää, mitkä positiot ovat kunkin position käytettävissä. Se tuo
laskentaan suunnan erottamalla aiemmat ja myöhemmät positiot, mutta
huomiopisteen sisältöperusteinen osa
\(q_i^{\top}k_j\) ei edelleenkään eksplisiittisesti riipu position
\(i\) absoluuttisesta sijainnista tai positioiden \(i\) ja \(j\)
välisestä etäisyydestä. Tämän vuoksi Transformer-malliin liitetään
positiokoodaus (engl. \emph{position embedding}), jonka avulla tokenien järjestys voidaan ottaa
huomioon niiden esityksiä tai huomiopisteitä muodostettaessa.

\subsubsection{Absoluuttinen positiokoodaus}

Absoluuttisessa positiokoodauksessa jokaiselle positiolle
\(i \in \{1,\dots,T\}\) muodostetaan vektori
\(p_i \in \mathbb{R}^{d}\), joka yhdistetään position tokenin
upotusvektoriin. Jos position \(i\) tokenin upotusvektori on
\(e_i \in \mathbb{R}^{d}\), Transformer-kerrosten syötteenä käytetään
esitystä

$$
h_i^{(0)} = e_i + p_i.
$$

Matriisimuodossa tämä voidaan kirjoittaa

$$
H^{(0)} = E + P,
$$

missä \(E \in \mathbb{R}^{T \times d}\) sisältää tokenupotukset ja
\(P \in \mathbb{R}^{T \times d}\) niitä vastaavat positiovektorit.
Koska molemmat vektorit kuuluvat samaan avaruuteen, niiden summa voidaan
antaa suoraan itsehuomiokerrokselle ilman esityksen dimension
kasvattamista.

Positiovektorit voidaan käsitellä opittavina parametreina. Tällöin
jokaiselle positiolle \(i \leq L\) opitaan oma vektori \(p_i\), missä
\(L\) on mallin suurin sallittu kontekstinpituus. Menetelmä on
yksinkertainen, mutta sitoo positioparametrit ennalta määrättyihin
positioihin eikä määrittele esityksiä positioille \(i > L\).

Alkuperäisessä Transformerissa käytettiin parametrien sijaan kiinteää
sinimuotoista positiokoodausta \cite{vaswani2017attention}. Sen
koordinaatit määritellään kaavoilla

$$
\begin{aligned}
  \operatorname{PE}(i,2k)
  &=
  \sin\left(
    \frac{i}{10000^{2k/d}}
  \right), \\
  \operatorname{PE}(i,2k+1)
  &=
  \cos\left(
    \frac{i}{10000^{2k/d}}
  \right),
\end{aligned}
$$

missä \(i\) on positio ja \(k\) indeksoi koordinaattipareja. Eri
koordinaatit käyttävät eri taajuuksia, joten koodaus esittää position
samanaikaisesti usealla mittakaavalla. Sinimuotoinen koodaus voidaan
laskea myös koulutuksessa käytetyn kontekstinpituuden ulkopuolisille
positioille. Tämä ei kuitenkaan yksin takaa, että malli kykenisi
hyödyntämään koulutuksessa kohtaamiaan pidempiä sekvenssejä.

\subsubsection{Suhteellinen positiokoodaus}

Kielellisten riippuvuuksien kannalta kahden tokenin suhteellinen etäisyys
on usein niiden absoluuttisia positioita olennaisempi. Suhteellisissa
positiokoodauksissa huomiomekanismiin lisätään tämän vuoksi positioiden
erotuksesta riippuva termi. Yksinkertainen yleinen muoto
huomiopisteelle on

$$
S_{ij}
=
\frac{q_i^{\top}k_j}{\sqrt{d_k}}
+
b(i-j)
+
M_{ij},
$$

missä \(b(i-j)\) on suhteellisesta etäisyydestä riippuva
positiovinouma ja \(M_{ij}\) on kausaalinen maski. Vinouma voi olla
opittava parametri, ennalta määrätty funktio tai positioiden etäisyyksiä
väleihin ryhmittelevä esitys \cite{shaw2018self}. Tällöin samoilla
suhteellisilla etäisyyksillä voidaan käyttää samoja parametreja
riippumatta siitä, missä kohdassa sekvenssiä tokenit esiintyvät.

Esimerkiksi ALiBi-menetelmässä kunkin huomiopään pisteisiin lisätään
etäisyyden mukana pienenevä lineaarinen vinouma
\cite{press2022train}. Huomiopään \(r\) pisteet voidaan kirjoittaa
muodossa

$$
S_{ij}^{(r)}
=
\frac{
  (q_i^{(r)})^{\top}k_j^{(r)}
}{\sqrt{d_k}}
-
m_r(i-j)
+
M_{ij},
\qquad j \leq i,
$$

missä \(m_r > 0\) on huomiopäälle määrätty kulmakerroin. Tällöin
kaukaisempien tokenien huomiopisteitä pienennetään enemmän kuin lähellä
olevien tokenien pisteitä.

\subsubsection{Kiertävä positiokoodaus}

Monissa moderneissa dekooderipohjaisissa kielimalleissa käytetään
kiertävää positiokoodausta (engl.\ \emph{rotary position embedding},
RoPE) \cite{su2021roformer}. Toisin kuin absoluuttisessa
positiokoodauksessa, RoPE ei lisää positiovektoria tokenin
syöte-esitykseen. Sen sijaan position \(i\) tieto liitetään
huomiomekanismiin kiertämällä kysely- ja avainvektorien
koordinaattipareja position mukaan.

Kahdessa ulottuvuudessa kulmaa \(i\theta\) vastaava kiertomatriisi on

$$
R(i,\theta)
=
\begin{pmatrix}
  \cos(i\theta) & -\sin(i\theta) \\
  \sin(i\theta) & \cos(i\theta)
\end{pmatrix}.
$$

Kysely- ja avainvektorien koordinaatit jaetaan pareihin, joille
käytetään eri taajuuksia. Kun \(d_k\) on parillinen, voidaan määritellä
lohkodiagonaalinen matriisi

$$
R_i
=
\operatorname{diag}\left(
  R(i,\theta_0),
  R(i,\theta_1),
  \dots,
  R(i,\theta_{d_k/2-1})
\right),
$$

missä taajuudet voidaan valita esimerkiksi muodossa

$$
\theta_k = 10000^{-2k/d_k}.
$$

Position \(i\) kysely ja position \(j\) avain muunnetaan tällöin
muotoon

$$
\widetilde{q}_i = R_i q_i,
\qquad
\widetilde{k}_j = R_j k_j.
$$

Kiertojen jälkeen laskettava huomiopiste on

$$
\begin{aligned}
  \widetilde{q}_i^{\top}\widetilde{k}_j
  &=
  (R_i q_i)^{\top}(R_j k_j) \\
  &=
  q_i^{\top}R_i^{\top}R_j k_j.
\end{aligned}
$$

Kiertomatriiseille pätee \(R_i^{\top}R_j = R_{j-i}\), joten

$$
\widetilde{q}_i^{\top}\widetilde{k}_j
=
q_i^{\top}R_{j-i}k_j.
$$

Huomiopisteen positioriippuvuus määräytyy siten positioiden erotuksen
\(j-i\) perusteella. RoPE käyttää siis kyselyiden ja avainten
absoluuttisia positioita kiertokulmien määrittämiseen, mutta niiden
pistetuloon syntyy suhteellisesta positiosta riippuva rakenne. Tämä
yhdistää absoluuttisen ja suhteellisen positiokoodauksen piirteitä ilman
erillisiä opittavia positiovektoreita.

Kiertoa sovelletaan tavallisesti kysely- ja avainvektoreihin mutta ei
arvovektoreihin. Kausaalisen monipäisen itsehuomion pää \(r\) voidaan
tällöin kirjoittaa muodossa

$$
\operatorname{head}_r
=
\operatorname{softmax}\left(
  \frac{
    \widetilde{Q}_r\widetilde{K}_r^{\top}
  }{\sqrt{d_k}}
  +
  M
\right)V_r,
$$

missä matriisien \(\widetilde{Q}_r\) ja
\(\widetilde{K}_r\) jokainen rivi on kierretty sitä vastaavan position
mukaan.

\subsubsection{Positiokoodaus ja kontekstinpituus}

Positiokoodaus vaikuttaa siihen, kuinka malli toimii koulutuksessa
käytettyä pidemmillä sekvensseillä. Opittavat absoluuttiset
positioupotukset eivät sellaisinaan määrittele esityksiä koulutetun
positioalueen ulkopuolelle. Sinimuotoiset koodaukset ja RoPE voidaan
puolestaan laskea mielivaltaisille positioille, mutta mallin ei silti
voida olettaa automaattisesti yleistävän huomattavasti pidempiin
sekvensseihin. Malli on koulutuksen aikana oppinut hyödyntämään vain
tietyn pituisilla sekvensseillä esiintyviä positioita ja etäisyyksiä.

Tämän vuoksi mallin arkkitehtuurin sallima kontekstiraja, koulutuksessa
käytetty kontekstinpituus ja mallin toiminnallinen kontekstinpituus on
hyödyllistä erottaa toisistaan. Arkkitehtuurin kontekstiraja ilmaisee
suurimman teknisesti sallitun syötepituuden, kun taas toiminnallinen
kontekstinpituus kuvaa sitä, kuinka pitkältä etäisyydeltä malli kykenee
käytännössä hyödyntämään informaatiota luotettavasti. Pitkän kontekstin
malleissa RoPE-kulmia voidaan lisäksi skaalata tai interpoloida, jotta
malli voidaan sovittaa alkuperäistä koulutuskontekstia pidemmille
sekvensseille.

\subsection{Dekooderikerros}

Yksittäinen dekooderikerros koostuu kausaalisesta monipäisestä itsehuomiosta ja
positiokohtaisesta myötäkytkentäverkosta. Kumpaankin komponenttiin liitetään
jäännöskytkentä (engl.\ \emph{residual connection}) sekä normalisointi, jotka
varmistavat laskennan numeerisen stabiilisuuden ja sujuvan gradienttien kulun
syvissä arkkitehtuureissa \cite{he2016deep, ba2016layer, vaswani2017attention}.

Laskennan vakauttamiseksi esitykset normalisoidaan ennen kuhunkin
alilohkoon syöttämistä. Normalisointifunktio \(\operatorname{Norm} \colon
\mathbb{R}^{d} \to \mathbb{R}^{d}\) sovelletaan itsenäisesti kunkin position
esitysvektoriin. Yleinen valinta nykyaikaisissa malleissa on neliöllisen
keskiarvon normalisointi (engl.\ \emph{root mean square normalization},
RMSNorm) \cite{zhang2019root}. Tämä skaalaa vektorin sen komponenttien neliöllisen keskiarvon avulla
$$
\operatorname{RMSNorm}(x) = \frac{x}{\operatorname{RMS}(x)} \odot \gamma,
$$
missä \(\odot\) on Hadamard-tulo, \(\gamma \in \mathbb{R}^d\)
on opittava skaalausparametri ja
$$
\operatorname{RMS}(x) = \sqrt{\frac{1}{d} \sum_{j=1}^{d} x_j^2 + \epsilon},
$$
missä pieni vakio \(\epsilon > 0\) estää nollalla jakamisen. Matriisille
\(H \in \mathbb{R}^{T \times d}\) normalisointi \(\operatorname{Norm}(H)\)
määritellään soveltamalla kuvausta riveittäin jokaiseen esitysvektoriin.

Itsehuomiokerrosta seuraa positiokohtainen myötäkytkentäverkko
(engl.\ \emph{feed-forward network}, FFN). Se vastaa luvussa \ref{sec:mlp}
tarkasteltua monikerroksista perseptronia, jota sovelletaan itsenäisesti ja
identtisesti kunkin position esitykseen. Kaksi affiinia muunnosta sisältävä
FFN voidaan määritellä funktiolle \(x \in \mathbb{R}^d\) muodossa
$$
\operatorname{FFN}(x) = \phi(x W_1 + b_1) W_2 + b_2,
$$
missä \(W_1 \in \mathbb{R}^{d \times d_{ff}}\), \(W_2 \in \mathbb{R}^{d_{ff} \times d}\),
\(b_1 \in \mathbb{R}^{d_{ff}}\), \(b_2 \in \mathbb{R}^d\) ja \(\phi\) on epälineaarinen
aktivaatiofunktio. Sisäisen piiloesityksen dimensio \(d_{ff}\) valitaan tyypillisesti
suuremmaksi kuin mallin dimensio \(d\) (usein \(d_{ff} = 4d\)). Matriisimuodossa
koko esitysmatriisille \(H \in \mathbb{R}^{T \times d}\) laskenta kirjoitetaan
vastaavasti muodossa \(\operatorname{FFN}(H) = \phi(H W_1 + \mathbf{1} b_1^{\top}) W_2 + \mathbf{1} b_2^{\top}\).

Koko dekooderilohko voidaan määritellä funktiona
\(\operatorname{DecoderBlock} \colon \mathbb{R}^{T \times d} \to \mathbb{R}^{T \times d}\).
Olkoon syöte esitysmatriisi \(H \in \mathbb{R}^{T \times d}\). Lohkon laskenta
etenee kahdessa vaiheessa:
$$
H^{(1)} = H + \operatorname{CausalMultiHead}(\operatorname{Norm}_1(H)),
$$
$$
\operatorname{DecoderBlock}(H) = H^{(1)} + \operatorname{FFN}(\operatorname{Norm}_2(H^{(1)})),
$$
missä \(\operatorname{Norm}_1\) ja \(\operatorname{Norm}_2\) ovat toisistaan
erillisiä normalisointeja omilla opittavilla parametreillaan. Jäännöskytkennän
ansioista syöte \(H\) summautuu suoraan alilohkojen muunnosten tulokseen, mikä
luo suoran gradienttiväylän verkon läpi ja mahdollistaa kymmenien tai satojen
kerrosten tehokkaan kouluttamisen ilman gradientin häviämistä.

Lohkon kahdella pääkomponentilla on toisiaan täydentävät tehtävät informaation käsittelyssä. Kausaalinen itsehuomio mahdollistaa sekvenssin eri positioiden välisen vuorovaikutuksen ja informaation reitittämisen tokenien välillä, kun taas positiokohtainen myötäkytkentäverkko vastaa kunkin position esitysvektorin riippumattomasta epälineaarisesta muunnoksesta ilman positioiden välistä tiedonsiirtoa.

\subsection{Seuraavan tokenin ennustaminen ja kielimallin tavoitefunktio}

Kuten luvussa \ref{sec:n-gram} esitetyn \(n\)-gram-mallin ja edellisessä
luvussa käsitellyn MLP-kielimallin tapauksessa, myös Transformer-pohjaisen
kielimallin tavoitteena on muodostaa ehdollinen todennäköisyysjakauma
seuraavalle tokenille. Erona on se, että Transformer muodostaa kontekstin
esityksen useiden itsehuomio- ja eteenpäinsyöttävien kerrosten avulla.

Olkoon \(h_t^{(L)}\in\mathbb{R}^{d}\) position \(t\) viimeisen
Transformer-kerroksen esitys. Sanaston tokenien normalisoimattomat
pistemäärät eli logitit määritellään muodossa
\begin{align*}
  z_t
  =
  W_{\mathrm{vocab}}h_t^{(L)}
  +
  b_{\mathrm{vocab}},
\end{align*}
missä
\begin{align*}
  W_{\mathrm{vocab}}
  &\in \mathbb{R}^{|V|\times d}, \\
  b_{\mathrm{vocab}}
  &\in \mathbb{R}^{|V|}.
\end{align*}
Logitvektorin \(z_t\) komponentti \(z_{t,j}\) vastaa tokenin \(v_j\)
pistemäärää. Seuraavan tokenin ehdollinen todennäköisyysjakauma saadaan
softmax-funktiolla:
\begin{align*}
  p_{\theta}(x_{t+1}=v_j\mid x_1,\ldots,x_t)
  =
  \frac{\exp(z_{t,j})}
  {\sum_{k=1}^{|V|}\exp(z_{t,k})}.
\end{align*}

Koko tokenisekvenssin todennäköisyys voidaan kirjoittaa
autoregressiivisena tulona:
\begin{align*}
  p_{\theta}(x_1,\ldots,x_T)
  =
  \prod_{t=1}^{T-1}
  p_{\theta}(x_{t+1}\mid x_1,\ldots,x_t).
\end{align*}
Käytännössä ensimmäisen tokenin ennustaminen voidaan määritellä
aloitustokenin avulla. Jos \(x_0=\langle s\rangle\), voidaan sama esitys
kirjoittaa muodossa
\begin{align*}
  p_{\theta}(x_1,\ldots,x_T\mid x_0)
  =
  \prod_{t=0}^{T-1}
  p_{\theta}(x_{t+1}\mid x_0,\ldots,x_t).
\end{align*}

Tämä on suora yleistys luvussa \ref{sec:n-gram} esitetystä
Markov-oletukseen perustuvasta kielimallista. \(n\)-gram-mallissa konteksti
rajoitetaan eksplisiittisesti viimeiseen \(n-1\) tokeniin, kun taas
Transformer-mallissa ehdollinen jakauma voi riippua kaikista
käytettävissä olevista aiemmista tokeneista. Käytännössä riippuvuutta
rajoittaa mallin suurin sallittu kontekstipituus.

Mallin parametrit \(\theta\) estimoidaan minimoimalla negatiivinen
log-uskottavuus eli tokenitason ristientropiahäviö:
\begin{align*}
  \mathcal{L}(\theta;X)
  =
  -
  \sum_{t=1}^{T-1}
  \log
  p_{\theta}
  \left(
    x_{t+1}
    \mid
    x_1,\ldots,x_t
  \right).
\end{align*}
Sekvenssin pituudella normalisoitu häviö on
\begin{align*}
  \overline{\mathcal{L}}(\theta;X)
  =
  -
  \frac{1}{T-1}
  \sum_{t=1}^{T-1}
  \log
  p_{\theta}
  \left(
    x_{t+1}
    \mid
    x_1,\ldots,x_t
  \right).
\end{align*}
Tavoitefunktio on siten periaatteellisesti sama kuin aiemmin käsitellyissä
maksimiuskottavuuteen perustuvissa kielimalleissa. Keskeinen ero on mallin
parametrisoinnissa: Transformer ei talleta erillistä jakaumaa jokaiselle
diskreetille kontekstille, vaan muodostaa ehdollisen jakauman syvän,
jaettuihin parametreihin perustuvan funktion avulla.

Koulutusvaiheessa koko syötesekvenssi voidaan käsitellä samanaikaisesti.
Kausaalinen maski varmistaa, ettei tokenia \(x_{t+1}\) käytetä sitä ennustavan
position \(t\) laskennassa. Tämän ansiosta kaikkien positioiden
ristientropiahäviöt voidaan laskea rinnakkain, vaikka mallin määrittelemä
todennäköisyysjakauma on autoregressiivinen. Parametrien gradientit saadaan
vastavirta-algoritmilla, ja parametreja päivitetään luvussa 4 esitetyillä
gradienttipohjaisilla menetelmillä. Käytännössä suurten kielimallien
esikoulutuksessa käytetään usein AdamW-optimointimenetelmää.

Tokenitason keskimääräisestä negatiivisesta log-uskottavuudesta voidaan
määritellä myös perplexiteetti:
\begin{align*}
  \operatorname{PPL}
  =
  \exp
  \left(
    \overline{\mathcal{L}}(\theta;X)
  \right).
\end{align*}
Perplexiteetti mittaa mallin keskimääräistä epävarmuutta havaittuja tokeneita
kohtaan. Pieni perplexiteetti tarkoittaa, että malli antaa aineiston
tokeneille keskimäärin suuren todennäköisyyden. Mittari ei kuitenkaan
sellaisenaan arvioi esimerkiksi vastauksen faktuaalista oikeellisuutta,
turvallisuutta, hyödyllisyyttä tai kykyä hyödyntää ulkoisia lähteitä.

\subsection{Autoregressiivinen generointi}

Koulutuksen jälkeen mallia käytetään tekstin generointiin iteratiivisesti.
Olkoon alkuperäinen syötekonteksti
\begin{align*}
  (x_1,\ldots,x_t).
\end{align*}
Malli muodostaa ensin jakauman
\begin{align*}
  p_{\theta}(x_{t+1}\mid x_1,\ldots,x_t).
\end{align*}
Tämän jälkeen jakaumasta valitaan uusi token \(\widehat{x}_{t+1}\), joka
liitetään kontekstiin. Seuraavalla askeleella malli muodostaa jakauman
\begin{align*}
  p_{\theta}
  \left(
    x_{t+2}
    \mid
    x_1,\ldots,x_t,\widehat{x}_{t+1}
  \right).
\end{align*}
Prosessia jatketaan, kunnes generoidaan lopetustokeni tai saavutetaan ennalta
määritelty enimmäispituus.

Yksinkertaisin valintasääntö on ahne dekoodaus, jossa valitaan aina
todennäköisin tokeni:
\begin{align*}
  \widehat{x}_{t+1}
  =
  \arg\max_{v_j\in V}
  p_{\theta}
  \left(
    x_{t+1}=v_j
    \mid
    x_1,\ldots,x_t
  \right).
\end{align*}
Vaihtoehtoisesti tokeni voidaan valita satunnaisotannalla mallin
ennustamasta jakaumasta. Tällöin generoinnin satunnaisuutta voidaan säätää
lämpötilaparametrilla \(\tau>0\):
\begin{align*}
  p_{\theta,\tau}(x_{t+1}=v_j\mid x_1,\ldots,x_t)
  =
  \frac{\exp(z_{t,j}/\tau)}
  {\sum_{k=1}^{|V|}\exp(z_{t,k}/\tau)}.
\end{align*}
Kun \(\tau<1\), jakauma keskittyy suuremman todennäköisyyden tokeneihin.
Kun \(\tau>1\), jakauma tasoittuu ja generointi muuttuu satunnaisemmaksi.
Käytännössä voidaan käyttää myös esimerkiksi top-\(k\)- ja top-\(p\)-
näytteenottoa, joissa näytteenotto rajoitetaan todennäköisimmiksi arvioituihin
tokeneihin.

Autoregressiivinen generointi poikkeaa koulutuksesta laskennallisesti siinä,
että uusi token lisätään kontekstiin yksi kerrallaan. Tehokkuuden
parantamiseksi toteutuksissa talletetaan aiemmissa generointivaiheissa
lasketut avain- ja arvovektorit niin sanottuun avain--arvovälimuistiin
(engl.\ \emph{key--value cache}). Tällöin aikaisempien tokenien avain- ja
arvoprojektioita ei tarvitse laskea uudelleen jokaisella generointiaskeleella.

\subsection{Skaalaus ja jälkikoulutus}

Suurten kielimallien suorituskykyyn vaikuttavat muun muassa parametrien määrä,
harjoitusaineiston määrä ja laatu, koulutuksessa käytettyjen tokenien määrä,
käytettävissä oleva laskentakapasiteetti sekä konteksti-ikkunan pituus.
Skaalauslakeja koskevissa tutkimuksissa on havaittu, että
kielimallien ennustehäviö yleensä pienenee säännönmukaisesti mallin,
aineiston ja laskennan määrän kasvaessa
\cite{kaplan2020scaling,hoffmann2022training}. Pelkkä parametrimäärä ei
kuitenkaan yksin määritä mallin laatua, sillä myös koulutusaineiston
koostumus, optimointi, arkkitehtuurivalinnat ja jälkikoulutus vaikuttavat
olennaisesti mallin toimintaan.

Esikoulutus tarkoittaa tavallisesti edellä määritellyn
seuraavan-tokenin-ennustustehtävän optimointia laajalla tekstiaineistolla.
Esikoulutettu malli oppii tekstin tilastollisia, syntaktisia ja semanttisia
säännönmukaisuuksia, mutta sen käyttäytyminen ei vielä välttämättä vastaa
vuorovaikutteisen avustajan käyttöä. Tämän vuoksi esikoulutusta seuraa usein
jälkikoulutus, jossa mallia sovitetaan ohjeita ja keskustelumuotoista
vuorovaikutusta sisältävään aineistoon.

Ohjatulla hienosäädöllä (engl.\ \emph{supervised fine-tuning}) malli
optimoidaan esimerkiksi käyttäjän ohjeen ja tavoitellun vastauksen muodostamilla
pareilla. Lisäksi mallin tuottamia vastauksia voidaan optimoida
ihmisten tai muiden arvioijien ilmaisemien preferenssien perusteella.
Yksi tunnettu lähestymistapa on ihmispalautteeseen perustuva
vahvistusoppiminen (engl.\ \emph{reinforcement learning from human feedback},
RLHF) \cite{ouyang2022training}. Jälkikoulutus ei muuta sitä perustavaa
tosiasiaa, että malli tuottaa tekstin autoregressiivisesti seuraavaa tokenia
ennustamalla. Se kuitenkin vaikuttaa merkittävästi siihen, millaisia
vastauksia malli todennäköisesti tuottaa käyttäjän antamissa tilanteissa.

\subsection{Parametrisen kielimallin rajoitteet ja yhteys RAG-järjestelmiin}

Transformer-pohjainen suuri kielimalli voi tuottaa sujuvaa ja
kontekstuaalisesti johdonmukaista tekstiä, mutta sen tietämys ei ole
tietokantamuodossa eikä sen parametreihin tallentunut informaatio ole
välttämättä ajantasaista, täydellistä tai helposti todennettavaa. Malli voi
myös tuottaa uskottavan tuntuisia väitteitä, jotka eivät perustu luotettavaan
lähteeseen tai ovat virheellisiä. Tätä ilmiötä kutsutaan usein
hallusinoinniksi, vaikka kyse ei ole mallin sisäisestä havaintokokemuksesta,
vaan todennäköisyyspohjaisen generointiprosessin tuottamasta virheellisestä tai
perusteettomasta tekstistä.

Lisäksi mallin käytettävissä oleva konteksti on rajallinen, ja mallin kyky
hyödyntää sille annettua kontekstia riippuu sekä syötteen rakenteesta että
mallin koulutuksesta. Se, että asiakirja sisällytetään mallin kontekstiin, ei
takaa, että malli käyttää asiakirjaa oikein tai että sen tuottama vastaus
perustuu yksinomaan kyseiseen asiakirjaan.

Hakua hyödyntävä generointi (engl.\ \emph{retrieval-augmented generation},
RAG) pyrkii lieventämään näitä rajoitteita liittämällä mallin
generointikontekstiin ulkoisesta tietolähteestä haettuja dokumentteja
\cite{lewis2020retrieval}. Tällöin kielimallin parametreihin sisältyvää
tietoa täydennetään generointihetkellä saatavilla olevalla dokumenttipohjaisella
tiedolla. Menetelmä voi parantaa erityisesti ajantasaisen, organisaatio- tai
toimialakohtaisen ja lähteistettävän tiedon hyödyntämistä. Se tuo kuitenkin
mukanaan uusia mahdollisia virhelähteitä, kuten epäonnistuneen haun,
epärelevanttien dokumenttien sisällyttämisen kontekstiin sekä dokumenttien
virheellisen tulkinnan.

Seuraavissa luvuissa tarkastellaan retrieval-augmented generation -järjestelmän
toteutusta ja arvioidaan, miten Transformer-pohjainen suuri kielimalli toimii
käytännön sovelluksessa, jossa generointia tuetaan ulkoisesta
tietolähteestä haetulla kontekstilla.
